% !TEX root = ../main.tex
\section{Data and Method}
\label{sec:method}

This section describes data collection, variable operationalization, and the
analytic strategies used to address each research question.

%% ----------------------------------------------------------------
\subsection{Data Collection}
\label{sec:data}

Our data come from the Apify Store's public REST API, which exposes actor-level
metadata, usage statistics, pricing structures, and developer information. The
API returns paginated JSON responses that include storefront fields plus
structured metadata such as event-level pricing and category identifiers. We
additionally query the actor-detail endpoint \texttt{GET /v2/acts/\{actorId\}}
to retrieve each tool's \texttt{actorPermissionLevel} and
\texttt{actorStandby}---the two architectural fields that operationalize
agentic readiness (\autoref{sec:variables}).

\paragraph{Enumeration and coverage.}
The global search endpoint truncates results at roughly 15{,}000 actors, so we
enumerate via the public sitemap: we extract 5{,}400+ developer usernames, then
issue per-developer paginated queries. We set \texttt{includeUnrunnableActors=true}
on every query so actors excluded by the API's default runnability filter are
retained. This two-stage strategy covers the full marketplace: every actor
belongs to exactly one developer, and every developer with a public actor
appears in the sitemap.

\paragraph{Snapshots and validation.}
We collected 17 daily snapshots (2026-08-19 to 2026-09-04); all
cross-sectional analyses use the latest (2026-09-04). The actor-detail endpoint
was queried on 2026-08-27, so the two architectural fields that define agentic
readiness (permission level and standby) are measured once. The remaining
snapshots serve for validation: zero duplicate rows, high actor retention, and
stability of the store's agentic-payments whitelist flag (recorded daily from
the store listing; only 1.4\% of actors change state across the 17 days).
Because the whitelist flag is a platform-side summary of agentic eligibility
that incorporates the architectural criteria among its inputs, its stability is
suggestive---but not conclusive---of architectural stability, which we cannot
observe daily. The single-snapshot census is 67{,}609 actors from
5{,}789 developers; among the 57{,}032 pay-per-event (PPE) actors, the detail
endpoint returned permission and standby fields for 53{,}502, of which 50{,}917
(95\%) are agentic-ready.

\paragraph{Ethics.} All data are public and require no authentication; the
dataset contains no personal data beyond public developer usernames.

\paragraph{Analytic samples.} Three nested samples structure the analyses: the
\emph{full marketplace} ($N = 67{,}609$), the \emph{PPE sample}
($N = 57{,}032$; 53{,}502 with architecture fields), and the
\emph{mixed-developer PPE sample} ($N = 343$ developers, 16{,}537 actors) that
publish both agentic-ready and non-agentic-ready tools and support the
within-developer fixed-effects design.

%% ----------------------------------------------------------------
\subsection{Variables and Operationalization}
\label{sec:variables}

\paragraph{Outcome.}
Our primary outcome is $\log(1 + \texttt{monthly\_users})$, the breadth of an
actor's user base over the trailing 30 days. We use the log-plus-one transform
because adoption counts are heavily right-skewed (median 1, top actors in the
tens of thousands); percentages reported as $\exp(\hat\beta) - 1$ refer to the
$1 + \texttt{monthly\_users}$ scale, not monthly users itself. As a secondary
outcome we use $\log(1 + \texttt{runs\_30d})$.

\paragraph{Treatment variable.}
The treatment is \emph{agentic readiness}, a binary indicator equal to one when
a PPE tool is configured with \emph{limited permissions} and \emph{no standby
mode}. We operationalize it from two actor-detail fields:
\texttt{actorPermissionLevel} ($\texttt{LIMITED\_PERMISSIONS}$ vs.\
$\texttt{FULL\_PERMISSIONS}$) and \texttt{actorStandby}/\texttt{standbyUrl}
(standby enabled vs.\ not). A tool is agentic-ready when it is
limited-permissions and non-standby. These two fields are the
\emph{architectural} component of Apify's agentic-payment eligibility
(\autoref{sec:context}): they are necessary conditions that the developer sets
in the tool's architecture, distinct from the developer-level KYC and
approval-level curation gates that the platform applies on top of them. Agentic
readiness thus describes \emph{architectural} admission to the agent-accessible
set---the tools the platform's architecture renders suitable for AI agents to
discover, invoke, and pay for---not the platform's full approval decision and
not realized agent use; we do not observe whether any given agent invokes a
tool.

\paragraph{Control variables.}
The within-developer specifications control for tool age
($\log(1 + \texttt{age\_days})$), description length, and primary-category fixed
effects, in addition to the developer fixed effects. As a robustness check we
also add further listing-time covariates---log per-event price, log build count,
and the number of pricing events. We do not control for ratings or review
counts, which are post-treatment adoption outcomes rather than confounds. The
per-category regressions instead control for $\log(1 + \texttt{total\_builds})$
(developer effort), developer portfolio size, and the number of pricing events,
because within-developer treatment variation is too sparse in most individual
categories to estimate developer fixed effects reliably.

\paragraph{Category membership.}
Actors are multi-labeled: 84.4\% of PPE actors carry two or three categories.
We define an actor's \emph{primary category} as the first entry in its
platform-returned category list and use primary-category indicators for fixed
effects and per-category regressions; as a robustness check we re-estimate the
per-category results under a multi-hot (any-membership) assignment
(\autoref{sec:results:rq3}).

%% ----------------------------------------------------------------
\subsection{Conceptual Framework}
\label{sec:framework}

Agentic readiness can be associated with adoption through two pathways
(\autoref{fig:mechanism}). The first is \emph{agent-mediated}: readiness
determines which tools enter the feasible interaction set for machine agents,
potentially changing routing and discoverability, which changes the probability
of first adoption and, through popularity feedback, downstream concentration.
The second is \emph{developer-mediated}: the architectural choices that make a
tool agentic-ready---limited permissions, no standby---may be correlated with a
tool's task type, maturity, or target audience. A developer may configure its
mass-market tools with limited permissions while reserving full permission for
complex, enterprise workflows, so readiness is associated with adoption even
without any agent-side routing effect. Such tool-level selection cannot be
ruled out without observing agent traffic.

\begin{figure}[t]
  \centering
  \begin{tikzpicture}[
      box/.style={draw, rounded corners=2pt, align=center, font=\scriptsize, inner sep=4pt, minimum width=2.3cm},
      obs/.style={-{Stealth[length=2mm]}, thick},
      inf/.style={-{Stealth[length=2mm]}, thick, dashed},
    ]
    \node[box] (flag) at (0,0) {Agentic readiness\\(limited permissions, no standby)};
    \node[box] (feasible) at (-2.8,-1.3) {Feasible interaction set\\for machine agents};
    \node[box] (dev) at (2.8,-1.3) {Safer, more restricted\\tool architecture};
    \node[box] (routing) at (-2.8,-2.5) {Routing /\\discoverability};
    \node[box] (choice) at (2.8,-2.5) {Easier to adopt};
    \node[box] (adopt) at (0,-3.7) {First adoption\\(extensive margin)};
    \node[box] (conc) at (0,-5.0) {Downstream\\concentration};
    \draw[obs] (flag) -- (feasible);
    \draw[obs] (flag) -- (dev);
    \draw[inf] (feasible) -- (routing);
    \draw[obs] (dev) -- (choice);
    \draw[inf] (routing) -- (adopt);
    \draw[obs] (choice) -- (adopt);
    \draw[obs] (adopt) -- (conc);
  \end{tikzpicture}
  \caption{Conceptual framework: two pathways from agentic readiness to
    adoption. Solid arrows mark associations we observe; dashed arrows mark
    steps we infer but do not observe. The agent-mediated and developer-mediated
    pathways converge on first adoption and cannot be separated with
    cross-sectional data.}
  \label{fig:mechanism}
\end{figure}

We observe the endpoints of both pathways---agentic readiness and the adoption
outcomes---but not the interior: our data contain no agent-side activity, so the
routing and developer-expectation steps are inferred rather than observed. Our
estimates therefore describe the adoption associated with admission to the
agent-accessible set; attributing this association to agent routing
specifically requires agent-side data. The research questions map onto the
model: RQ1 estimates the total association, RQ2 locates the margin, and RQ3
characterizes concentration.

%% ----------------------------------------------------------------
\subsection{RQ1: Quantifying the Adoption Association of Agentic Readiness}
\label{sec:methods}

We estimate the association between agentic readiness and adoption with a
primary within-developer design and a set of robustness checks.

\subsubsection{Within-Developer Fixed Effects}

The primary specification exploits the 343 developers who publish \emph{both}
agentic-ready and non-agentic-ready pay-per-event tools (16{,}537 actors). For
these mixed developers we estimate
\begin{equation}
\label{eq:fe}
\log(1 + \text{MU}_{ij}) = \beta \cdot \text{Ready}_{ij}
  + \mathbf{X}_{ij} \boldsymbol{\gamma} + \alpha_j + \epsilon_{ij}
\end{equation}
where $\alpha_j$ is a developer fixed effect. The fixed effect absorbs every
time-invariant developer characteristic---reputation, marketing reach, audience
size, category specialization, pricing strategy---so $\beta$ identifies the
\emph{within-developer} contrast between a developer's agentic-ready and
non-agentic-ready tools. This is the strongest comparison our data support: it
holds the developer constant and asks whether readiness predicts adoption
across the tools a single developer publishes. The control vector
$\mathbf{X}_{ij}$ contains tool age, description length, and primary-category
fixed effects; standard errors are clustered at the developer level
\cite{wooldridge2010econometric}.

The design does not address \emph{within-developer selection}: a developer may
configure its most promising tools to be agentic-ready and invest more in them
\cite{boudreau2012complementors}. We treat this as a substantive alternative
explanation rather than a technical nuisance, and return to it in
\autoref{sec:discussion}.

\paragraph{Two estimands.} Developer size correlates with the readiness gap,
so we report two co-equal estimands. The \emph{equal-developer-weighted}
estimand down-weights observations by developer portfolio size: each tool of
developer $j$ receives weight $1/n_j$ (where $n_j$ is the developer's tool
count), so a developer's influence on the estimate does not grow with the number
of tools it publishes. This approximates asking whether the \emph{typical
developer's} agentic-ready tools outperform its non-agentic-ready tools; it
removes the size weighting but not the within-developer variance in readiness
share, so it is not an exact simple average of developer-level contrasts. The
\emph{actor-weighted} estimand weights each tool equally, asking whether the
\emph{typical tool} in the agentic-ready set outperforms the typical tool
outside it. The two answer different questions---one about developers, one about
tools---and we report both throughout (\autoref{sec:results:rq1}).

\paragraph{Who are the mixed developers?} The within-developer design rests on
the 343 mixed developers, a specific subsample rather than a random draw of the
marketplace (\autoref{tab:mixed_developers}). Relative to non-mixed developers,
they run larger portfolios (median 11 versus 2 tools), publish older tools
(median 105 versus 77 days), and attract far more users on average (mean 73.5
versus 11.9 monthly users). They are exclusively boutique-to-whale
developers---no single-tool developer can be mixed by construction. The
within-developer estimates therefore speak to the marketplace's prolific,
established developers rather than to its long tail, a boundary we carry into
the interpretation.

\begin{table}[t]
  \centering
  \caption{The 343 mixed developers versus non-mixed pay-per-event developers.
    Mixed developers publish both agentic-ready and non-agentic-ready tools and
    identify the within-developer effect; non-mixed developers publish only one
    type. Portfolio and age are developer-level; monthly users and agentic-ready
    share are tool-level.}
  \label{tab:mixed_developers}
  \small
  \begin{tabular}{lcc}
    \toprule
    & \textbf{Mixed} & \textbf{Non-mixed} \\
    \midrule
    Developers & 343 & 3{,}826 \\
    Median full-store portfolio & 11 & 2 \\
    Median tool age (days) & 105 & 77 \\
    Mean monthly users & 73.5 & 11.9 \\
    Agentic-ready share (\%) & 91.2 & 97.0 \\
    Solo / Boutique / Mid / Whale (\%) & 0 / 46 / 42 / 12 & 42 / 41 / 15 / 2 \\
    \bottomrule
  \end{tabular}
\end{table}


\subsubsection{Robustness}

Four checks bound the interpretation. First, we exclude \emph{whale developers}
($\geq 100$ pay-per-event tools) and re-estimate the within-developer model, to
rule out that a handful of large developers drive the result. Second, we use
propensity score matching (PSM) \cite{rosenbaum1983propensity, stuart2010matching},
estimating the propensity score from \emph{listing-time} covariates only (tool
age, description length, developer portfolio size, category) so post-treatment
signals cannot enter the model. Because the treated-to-control ratio is 19.7:1,
1:1 matching covers only a small fraction of agentic-ready tools; we therefore
report full-sample weighting estimators---overlap weighting (overlap
population, ATO) and inverse-probability-of-treatment weighting (treated
population, ATT)---alongside 1:1 matching. Because the comparison group is
small and potentially selected, we report covariate-balance diagnostics
(standardized mean differences before and after weighting) and the
propensity-score overlap across the two groups. Third, because
$\log(1+\texttt{monthly\_users})$ compresses the multiplicative gap near zero
(the median tool has one monthly user), we also fit a Poisson
pseudo-maximum-likelihood model with developer fixed effects
\cite{santossilva2006loggravity}, which estimates the gap on the monthly-user
count scale directly. Fourth, to address the
alternative explanation that agentic-ready and non-agentic-ready tools are
simply different tool types, we match each ready tool to its most similar
non-ready tool \emph{within the same developer} (1:1 nearest-neighbor, with
replacement, no caliper) under two similarity notions: description cosine
similarity from TF-IDF vectors (1--2 word n-grams, top 5{,}000 terms), and a
covariate distance (standardized Euclidean distance on log price, log age, and
description length, plus a bonus for an exact category match). For each matched
set we report the mean log-adoption gap with a developer-clustered bootstrap
95\% confidence interval---we resample developers (with replacement) and
recompute the mean gap over the pairs of the sampled developers, since matched
pairs are nested within developers and treating them as independent would
understate uncertainty---and the standardized covariate differences before and
after matching. Because nearest-neighbor matching is a nonsmooth estimator for
which the standard bootstrap is not generally valid
\cite{abadie2006largesample,abadie2008bias}, we treat these intervals as
descriptive rather than as exact inference for the average treatment effect on
the treated.

%% ----------------------------------------------------------------
\subsection{RQ2: Locating the Margin}
\label{sec:methods-rq2}

RQ2 asks \emph{where} the adoption association operates. We decompose it into
an \emph{extensive} margin---whether a tool has any monthly users at all---and
an \emph{intensive} margin---how many users it has conditional on having any.
The distinction maps onto the two pathways in \autoref{fig:mechanism}: an
extensive-margin association is consistent with readiness changing whether a
tool is ever discovered and first adopted, while an intensive-margin
association is consistent with readiness changing how far a tool scales once
adopted. Because the two pathways make different predictions about
\emph{where} the association appears, locating the margin is descriptive.

We estimate three outcomes within the mixed-developer, equal-developer-weighted
design: a linear-probability model for having any monthly users (extensive
margin), an OLS model on $\log(1 + \texttt{monthly\_users})$ conditional on
traction (intensive margin), and a linear-probability model for entering the
top decile of monthly users, which captures the extreme upper tail of adoption.
Each uses developer fixed effects, developer-clustered standard errors, and the
same controls as \autoref{eq:fe}.

%% ----------------------------------------------------------------
\subsection{RQ3: Heterogeneity and Concentration}
\label{sec:methods-rq3}

RQ3 asks two questions: whether the adoption association is uniform across
marketplace categories, and whether agentic readiness is associated with how
concentrated adoption is.

\subsubsection{Category Heterogeneity}

We re-estimate the adoption association separately within each category that
has at least 200 tools and at least 30 in each readiness group (15 categories),
using developer-clustered standard errors and controls for developer effort
($\log(1 + \texttt{total\_builds})$), developer portfolio size, and the number
of pricing events. Because actors are multi-labeled, we report both a
\emph{primary-category} assignment (the first category in the platform-returned
list) and a \emph{multi-hot} (any-membership) assignment, the latter to guard
against tag-ordering artifacts; we apply a Bonferroni correction within each
assignment. Because these per-category regressions use tool-level controls
rather than developer fixed effects, we also estimate a Ready $\times$ Category
interaction in the within-developer model and report within-developer,
within-category estimates where the design supports them, so the category
gradient is not inferred from between-developer variation alone.

\subsubsection{Concentration}

We measure concentration within the agentic-ready and non-agentic-ready sets
separately using three complementary statistics: the Gini coefficient (overall
inequality), top-1\%, top-5\%, and top-10\% user shares (the tail), and a Theil
decomposition into within- and between-category components. We also compare
within-category Gini coefficients across the two sets. The question is whether
admitting agents to a tool set is associated with a more winner-take-all
distribution of users.
